% Theme: Sorbet  --  best for intuition posts and beginner explainers
% Build:  ./build.sh example-sorbet.tex   (or compile on Overleaf)
\documentclass[12pt]{beamer}
\usetheme{igsorbet}

\ighandle{@quantum_danesh}
\igseries{Qubits desde cero}
\igissue{04}

\begin{document}

\coverslide{Intuición}
  {¿Por qué no se puede copiar un qubit?}
  {El teorema de no clonación, en cuatro pasos y sin miedo.}

\begin{frame}{El deseo}{Paso 1}
Queremos una máquina que copie \key{cualquier} estado:
\[
  U\ket{\psi}\ket{0} = \ket{\psi}\ket{\psi}
  \qquad \text{para todo } \ket{\psi}.
\]
En un sistema cerrado la evolución es unitaria, así que $U$ conserva los
productos escalares.
\end{frame}

\begin{frame}{Probamos con dos estados}{Paso 2}
\begin{align*}
  U\ket{\psi}\ket{0} &= \ket{\psi}\ket{\psi}\\
  U\ket{\phi}\ket{0} &= \ket{\phi}\ket{\phi}
\end{align*}
Producto escalar de ambas:
\[
  \braket{\psi}{\phi}\,\underbrace{\braket{0}{0}}_{=\,1}
  = \mhl<2->{\braket{\psi}{\phi}^{2}}
\]
\uncover<2->{A la izquierda queda $\braket{\psi}{\phi}$, a la derecha su cuadrado.}
\end{frame}

\begin{frame}{Solo dos opciones}{Paso 3}
Si un número cumple $x = x^2$, entonces $x = 0$ o $x = 1$:
\keyeq{\braket{\psi}{\phi} \in \{0,\,1\}}
Solo se pueden copiar estados \hl{idénticos u ortogonales}.
Un qubit desconocido, no.
\end{frame}

\begin{frame}{Copiar sería medir sin molestar}{Intuición}
\begin{notebox}
Si pudieras clonar, harías mil copias y medirías cada una. Conocerías el
estado sin perturbarlo, y eso la mecánica cuántica no lo permite.
\end{notebox}
\begin{defbox}[¿Y los bits clásicos?]
$\ket{0}$ y $\ket{1}$ son ortogonales, así que sí se copian:
$\mathrm{CNOT}\,\ket{x}\ket{0} = \ket{x}\ket{x}$.
\end{defbox}
\end{frame}

\begin{frame}{Por qué importa}{Paso 4}
\begin{itemize}
  \item \key{Criptografía cuántica}: un espía no puede copiar los qubits
        sin dejar huella.
  \item \key{Corrección de errores}: se reparte la información sin copiarla,
\end{itemize}
\[
  \alpha\ket{0} + \beta\ket{1} \;\longmapsto\;
  \alpha\ket{000} + \beta\ket{111} \;\neq\; \mhl*{\ket{\psi}^{\otimes 3}}
\]
{\footnotesize\color{igMuted}Wootters y Zurek, Nature \textbf{299}, 802 (1982);
Dieks, Phys.\ Lett.\ A \textbf{92}, 271 (1982).\par}
\end{frame}

\closingslide{Copiar, no. Compartir, sí.}
  {Guárdalo para cuando alguien te hable de teletransporte cuántico.}

\end{document}
